\section{Introducción: el puente entre el razonamiento matemático informal y el formal}

Esta clase la imparte Sean Welleck, de CMU, quien examina cómo aprovechar la IA para cerrar la brecha entre las \textbf{matemáticas informales} (informal math) y las \textbf{matemáticas formales} (formal math), combinando las ventajas de ambas.

\begin{importantbox}{Comparación entre las dos representaciones matemáticas}
\textbf{Matemáticas informales}: se expresan en lenguaje natural (texto, LaTeX, imágenes, etc.), son sumamente flexibles, pero no garantizan la corrección. El LLM puede cometer errores sutiles en los pasos intermedios.\\
\textbf{Matemáticas formales}: se escriben en un lenguaje de programación (Lean, Coq, Isabelle); que el código compile equivale a que la demostración sea correcta. Ofrecen una garantía absoluta de corrección, pero el umbral para escribirlas es alto y el proceso es laborioso.
\end{importantbox}

\subsection{Introducción a los asistentes de demostración formal (Lean)}

Lean es un asistente de demostración interactivo (Interactive Theorem Prover); sus conceptos centrales son:
\begin{itemize}
    \item \textbf{Enunciado del teorema}: expresar en lenguaje formal la proposición que se debe demostrar
    \item \textbf{Estado de la demostración} (Proof State): el contexto actual (las hipótesis) y el objetivo pendiente
    \item \textbf{Táctica} (Tactic): la regla de inferencia que se aplica en cada paso, la cual modifica el estado de la demostración
    \item \textbf{QED}: la demostración se completa cuando se resuelven todos los objetivos
\end{itemize}

\begin{knowledgebox}{Mathlib: la infraestructura de las matemáticas formales}
Mathlib es el proyecto de biblioteca matemática de Lean, construido en conjunto por numerosos colaboradores con las definiciones y los teoremas fundamentales de las matemáticas. Permite que las demostraciones formales reutilicen resultados ya existentes; el matemático Terence Tao ya ha logrado formalizar en Lean los teoremas de sus artículos.
\end{knowledgebox}

\subsection{La importancia de las matemáticas formales para la IA}

\begin{enumerate}
    \item \textbf{Salida verificable}: evita el razonamiento matemático incorrecto
    \item \textbf{Señal de recompensa perfecta}: se puede usar para el entrenamiento con RL, ya que si la demostración compila, es correcta
    \item \textbf{Banco de pruebas para el razonamiento}: abarca problemas de razonamiento desde sencillos hasta arbitrariamente difíciles
\end{enumerate}

La demostración formal de teoremas ha avanzado rápidamente en los últimos años: desde GPT-f (2020) hasta DeepSeek-Prover, los resultados en el banco de pruebas mini-F2F han subido de manera sostenida.
